% --- Содержимое файла materials/6_1.tex ---

\subsection{Теорема Рисса--Фреше}

\begin{definition}[Линейный функционал]
	Линейный оператор, действующий из нормированного пространства $X$ в скалярное поле $\mathbb{K}$ (где $\mathbb{K} = \mathbb{R}$ или $\mathbb{C}$), называется \textbf{линейным функционалом}.
	
	Множество всех линейных непрерывных функционалов над $X$ образует \textbf{сопряженное пространство}:
	\[ X^* = \mathcal{L}(X, \mathbb{K}) \]
\end{definition}

\begin{idea}[Смысл теоремы Рисса--Фреше]
	\textbf{Функционал = Вектор.}
	В «хороших» (гильбертовых) пространствах любой инструмент измерения (функционал) сам является вектором этого пространства.
	
	Измерение длины проекции вектора $x$ на вектор $y$ записывается как скалярное произведение $(y, x)$. Теорема утверждает, что \textbf{любой} линейный непрерывный функционал $f(x)$ выглядит именно так. Нет никаких других «экзотических» функционалов.
	
	Геометрически вектор $y$ — это градиент функционала, перпендикуляр к его ядру (нулевому уровню).
\end{idea}



\begin{theorem}[Рисса--Фреше]\label{Riss-Freshe}
	Пусть $H$ — гильбертово пространство, $f \in H^*$ — линейный ограниченный функционал.
	Тогда существует \textbf{единственный} вектор $y \in H$ такой, что функционал представим в виде скалярного произведения:
	\[ \forall x \in H: \quad f(x) = \langle y, x \rangle. \]
	При этом их нормы совпадают: $\|f\|_{H^*} = \|y\|_H$.
\end{theorem}

\begin{proof}
	\textbf{1. Существование.}
	
	\begin{itemize}
		
		\item \textit{Случай $f \equiv 0$.}
		
		Берём $y=0$. Тогда для любого $x\in H$:
		$$
		f(x)=0=\langle 0,x\rangle,
		$$
		и $\|f\|=\|y\|=0$.
		
		Далее считаем $f\not\equiv 0$.
		
		\medskip
		
		\item \textbf{Шаг 1. Рассматриваем ядро функционала.}
		
		Обозначим
		$$
		M:=\ker f=\{x\in H:\ f(x)=0\}.
		$$
		
		Тогда:
		\begin{itemize}
			\item $M\neq H$ (так как $f\neq 0$);
			\item $M$ замкнуто, поскольку $f$ непрерывен.
		\end{itemize}
		
		По теореме о разложении гильбертова пространства:
		$$
		H = M \oplus M^\perp,
		$$
		поэтому существует вектор
		$$
		b\in M^\perp,\quad b\neq 0.
		$$
		
		\medskip
		
		\item \textbf{Шаг 2. Стратегия.}
		
		Мы хотим получить формулу вида
		$$
		f(x)=C\langle b,x\rangle.
		$$
		
		Так как $f=0$ на $M$, естественно попытаться для каждого $x$
		вычесть из него компоненту вдоль $b$, чтобы попасть в ядро.
		
		То есть ищем число $\alpha$, для которого
		$$
		x-\alpha b\in M=\ker f.
		$$
		
		\medskip
		
		\item \textbf{Шаг 3. Подбор коэффициента.}
		
		Требуем:
		$$
		f(x-\alpha b)=0.
		$$
		
		По линейности:
		$$
		f(x-\alpha b)=f(x)-\alpha f(b).
		$$
		
		Так как $b\notin \ker f$ (иначе $b\perp b\Rightarrow b=0$),
		имеем $f(b)\neq 0$. Поэтому
		$$
		f(x)-\alpha f(b)=0
		\quad\Rightarrow\quad
		\alpha=\frac{f(x)}{f(b)}.
		$$
		
		Положим
		$$
		p_x:=x-\frac{f(x)}{f(b)}\,b.
		$$
		Тогда $p_x\in\ker f$.
		
		\medskip
		
		\item \textbf{Шаг 4. Используем ортогональность.}
		
		Так как $b\perp\ker f$, получаем:
		$$
		\langle b,p_x\rangle=0.
		$$
		
		Подставляя $p_x$:
		$$
		0=\left\langle b,\;x-\frac{f(x)}{f(b)}b\right\rangle
		=\langle b,x\rangle-\frac{f(x)}{f(b)}\|b\|^2.
		$$
		
		\medskip
		
		\item \textbf{Шаг 5. Выражаем функционал.}
		
		Из последнего равенства:
		$$
		f(x)=\frac{f(b)}{\|b\|^2}\,\langle b,x\rangle.
		$$
		
		Перенеся скаляр внутрь скалярного произведения (в первое место), получаем
		$$
		f(x)=\left\langle
		\frac{\overline{f(b)}}{\|b\|^2}\,b,\;x
		\right\rangle.
		$$
		
		Обозначим
		$$
		y:=\frac{\overline{f(b)}}{\|b\|^2}\,b.
		$$
		
		Тогда для всех $x\in H$:
		$$
		f(x)=\langle y,x\rangle.
		$$
		
		Существование доказано.
		
	\end{itemize}
	
	\medskip
	
	\textbf{2. Единственность.}
	
	Пусть
	$$
	\langle y,x\rangle=\langle z,x\rangle
	\quad \forall x\in H.
	$$
	
	Тогда
	$$
	\langle y-z,x\rangle=0\quad \forall x.
	$$
	
	Подставляя $x=y-z$, получаем:
	$$
	\|y-z\|^2=0 \;\Rightarrow\; y=z.
	$$
	
	\medskip
	
	\textbf{3. Равенство норм.}
	
	\begin{itemize}
		
		\item \textit{Оценка сверху.}
		
		По неравенству Коши–Буняковского:
		$$
		|f(x)|=|\langle y,x\rangle|
		\le \|y\|\,\|x\|.
		$$
		
		Следовательно,
		$$
		\|f\|=\sup_{\|x\|=1}|f(x)|
		\le \|y\|.
		$$
		
		\medskip
		
		\item \textit{Оценка снизу.}
		
		Если $y\neq 0$, берём $x=y$:
		$$
		f(y)=\langle y,y\rangle=\|y\|^2.
		$$
		
		По определению нормы функционала:
		$$
		|f(y)|\le \|f\|\,\|y\|.
		$$
		
		Отсюда:
		$$
		\|y\|^2\le \|f\|\,\|y\|
		\;\Rightarrow\;
		\|y\|\le \|f\|.
		$$
		
		(При $y=0$ равенство очевидно.)
		
	\end{itemize}
	
	Из двух неравенств:
	$$
	\|f\|=\|y\|.
	$$
	
\end{proof}
